% Derived LaTeX Rendering of IntelliTwin Authoritative Hypothesis Framework
% Primary Authority: research/hypotheses/hypothesis_framework.yaml & research/research_design/statistical_inference_protocol.md
% Pinned Source Commit: bb9e34579343881a831e3f052608ee15e183df76
% Governing Architecture Decision Record: ADR-018

\subsection*{Governance Rules on Hypothesis Integrity (ADR-018)}

\begin{quote}
\textbf{Authoritative Precedence:} This document is a derived LaTeX rendering. The YAML specification (\texttt{hypothesis\_framework.yaml}) and statistical protocol (\texttt{statistical\_inference\_protocol.md}) remain the primary immutable authorities.

\textbf{Decision Outcome Rules:} Each research question possesses a pre-specified primary outcome, comparator set, frozen numerical practical-effect threshold, and cluster-bootstrap decision rule classifying results into a three-way outcome: \textbf{Supported}, \textbf{Unsupported}, or \textbf{Inconclusive}.

\textbf{Outcome Hierarchy:} The primary outcome determines the confirmatory claim for each research question. Secondary outcomes verify operational constraints, reliability trade-offs, and failure safeguards. A favorable secondary outcome cannot replace failure on the primary outcome.

\textbf{No Post-Hoc Alteration:} Hypotheses, comparators, subgroups, thresholds, and decision rules are frozen and must not be altered after observing confirmatory test data. Any post-hoc modification must be explicitly designated as an exploratory analysis.

\textbf{Statistical Sample Unit:} The independent statistical unit is the physical \textbf{engine unit} ($i \in \{1, \dots, N\}$). Individual operational cycles ($t$) within an engine trajectory are autocorrelated and must \textbf{NEVER} be treated as independent statistical observations.
\end{quote}

\subsection*{Numerical Evaluation Contract (Protocol \texttt{step06-v2})}

\begin{itemize}
  \item \textbf{Primary Maintenance Cost Scenario}: \texttt{primary\_reference} ($\CPM = 1.0, \CFAIL = 5.0, \CWASTE = 0.20$ SCU per engine).
  \item \textbf{Secondary Sensitivity Scenarios}: \texttt{failure\_sensitive} ($\CFAIL = 10.0$) and \texttt{waste\_sensitive} ($\CWASTE = 0.40$).
  \item \textbf{Nominal Coverage}: Primary confirmatory coverage 90\% ($\alpha = 0.10$); descriptive sensitivities at 80\% and 95\%.
  \item \textbf{Cost Improvement Margin}: $\Deltacost > 0.20$ SCU per engine.
  \item \textbf{Coverage Error Improvement Margin}: $\Deltaerror > 0.01$ (one percentage point error reduction).
  \item \textbf{Maximum Absolute Coverage Error Ceiling}: $E < 0.05$ (prespecified 85\%--95\% empirical acceptance band).
  \item \textbf{Interval Noninferiority Margins}: $\Dwidth < 1.0$ cycle for MPIW and $\Dscore < 1.0$ cycle for Winkler interval score.
  \item \textbf{Robustness Non-Harm Margins}: $\Deltacost > -0.20$ SCU per engine and $\Deltaerror > -0.01$ in every cell.
  \item \textbf{Failure Safeguard}: Observed fleet failure rate increase $\Deltafail \le 0.0$ against comparator.
  \item \textbf{Fixed Model Seeds}: Five prespecified training seeds: \texttt{[2026, 2027, 2028, 2029, 2030]}.
\end{itemize}

\subsection*{Statistical Resampling and Multiplicity Control}

\begin{enumerate}
  \item \textbf{Resampling Engine}: 10,000 paired engine-cluster bootstrap resamples with replacement, stratified by dataset lifetime quartiles ($Q_1$--$Q_4$) fixed from split manifests.
  \item \textbf{PRNG Specification}: \texttt{numpy.random.Generator} backed by \texttt{PCG64} with seed sequence \texttt{SeedSequence([2026, 6, dataset\_number])}.
  \item \textbf{Decision Confidence Level}: Confirmatory decision rules use \textbf{98.75\% cluster-bootstrap percentile intervals} (quantiles 0.00625 and 0.99375), providing Bonferroni family-wise error rate control ($\alpha = 0.05 / 4 = 0.0125$) across the four primary research questions. Descriptive 95\% intervals (quantiles 0.025 and 0.975) are reported for context.
  \item \textbf{Subgroup Eligibility \& Bound Propagation}: Subgroup eligibility ($\ge 10$ engines) is frozen before resampling. If an eligible subgroup is absent in replicate $b$, its absolute error is bounded by $[0.0, 0.90]$ at 90\% nominal coverage, propagating conservatively.
  \item \textbf{Intersection-Union Testing (IUT)}: An RQ is Supported if and only if \textbf{every required component condition passes} at the 98.75\% decision interval level.
\end{enumerate}

\subsection*{Formal Hypothesis Definitions}

\subsubsection*{Research Question 1 (RQ1): Sequential Maintenance Utility and Coverage Reliability}
\begin{itemize}
  \item \textbf{Question}: \textit{Can a decision-aware, condition-adaptive conformal RUL calibration method improve sequential maintenance utility while preserving reliable prediction-interval coverage?}
  \item \textbf{Null Hypothesis ($H_{0,1}$)}: The proposed method does not improve sequential maintenance utility by at least 0.20 SCU per engine over both global and standard condition-aware calibration, or fails to maintain required coverage reliability.
  \item \textbf{Alternative Hypothesis ($H_{1,1}$)}: The proposed method improves maintenance utility by at least 0.20 SCU per engine over both global and standard condition-aware calibration while maintaining required coverage reliability.
  \item \textbf{Primary Outcome}: Total simulated maintenance cost per engine contrast ($\Deltacost = \overline{\text{Cost}}_B - \overline{\text{Cost}}_P$).
  \item \textbf{Primary Comparators}: Global conformal calibration (\texttt{global\_hcp}) and standard condition-aware calibration (\texttt{regime\_hcp}) on C-MAPSS FD002 held-out test split under \texttt{primary\_reference} at 90\% nominal coverage (both GBR and LSTM required).
  \item \textbf{Supported Logic}: For both GBR and LSTM, lower 98.75\% CI of $\Deltacost > 0.20$ SCU against both baselines, upper 98.75\% CI of marginal, worst-regime, and worst-stage coverage error $< 0.05$, and observed fleet failure rate increase $\le 0.0$ against both baselines.
  \item \textbf{Unsupported Logic}: For either model, upper 98.75\% CI of $\Deltacost < 0.20$ SCU against either baseline, lower 98.75\% CI of any coverage error $> 0.05$, or observed fleet failure rate increases over either baseline.
  \item \textbf{Inconclusive Logic}: No condition fails, but at least one decision interval covers its boundary, or inference is degenerate.
\end{itemize}

\subsubsection*{Research Question 2 (RQ2): Conditional Subgroup Reliability and Interval Sharpness}
\begin{itemize}
  \item \textbf{Question}: \textit{Can the proposed method improve reliability across operating regimes and degradation conditions without excessively widening prediction intervals?}
  \item \textbf{Null Hypothesis ($H_{0,2}$)}: The proposed method does not improve worst-group conditional coverage error by at least 0.01 over both baselines, or causes interval width or interval score to deteriorate by 1.0 cycle or more.
  \item \textbf{Alternative Hypothesis ($H_{1,2}$)}: The proposed method improves worst-group conditional coverage error by at least 0.01 over both baselines while restricting interval width and score deterioration to under 1.0 cycle.
  \item \textbf{Primary Outcome}: Worst-group conditional coverage error contrast ($\Deltaerror = \WorstErr_B - \WorstErr_P$).
  \item \textbf{Required Subgroup Families}: Operating regimes ($C_1$--$C_6$) and degradation lifecycle stages (Early, Mid, Late thirds).
  \item \textbf{Supported Logic}: For both models and both subgroup families, lower 98.75\% CI of $\Deltaerror > 0.01$ against both baselines, upper 98.75\% CI of worst-group error $< 0.05$, upper 98.75\% CI of $\Dwidth < 1.0$ cycle, and upper 98.75\% CI of $\Dscore < 1.0$ cycle.
  \item \textbf{Unsupported Logic}: For either model or subgroup family, upper 98.75\% CI of $\Deltaerror < 0.01$ against either baseline, lower 98.75\% CI of worst-group error $> 0.05$, lower 98.75\% CI of $\Dwidth > 1.0$ cycle, or lower 98.75\% CI of $\Dscore > 1.0$ cycle.
  \item \textbf{Inconclusive Logic}: No condition fails, but at least one decision interval covers its boundary, or inference is degenerate.
\end{itemize}

\subsubsection*{Research Question 3 (RQ3): Decision-Value Attribution vs. Matched Controls}
\begin{itemize}
  \item \textbf{Question}: \textit{Are the maintenance benefits of the proposed method attributable to decision-aware uncertainty calibration rather than simply more conservative prediction intervals or earlier maintenance?}
  \item \textbf{Null Hypothesis ($H_{0,3}$)}: The proposed method does not improve maintenance utility by at least 0.20 SCU per engine over every prespecified matched-conservatism control family, or violates coverage constraints.
  \item \textbf{Alternative Hypothesis ($H_{1,3}$)}: The proposed method provides maintenance improvements beyond those achievable through matched-conservatism baselines, demonstrating useful decision information rather than simple over-conservatism.
  \item \textbf{Primary Outcome}: Maintenance cost improvement over matched-conservatism controls ($\Deltacost = \overline{\text{Cost}}_{\text{control}} - \overline{\text{Cost}}_P$).
  \item \textbf{Four Evaluated Control Configurations across Three Families}:
    \begin{enumerate}
      \item \texttt{matched\_width\_interval}: $\lambda$-scaled baseline interval matching proposed validation MPIW within $\le 1.0$ cycle (GBR: \texttt{global\_hcp}, $\lambda=0.92$; LSTM: \texttt{regime\_hcp}, $\lambda=0.90$).
      \item \texttt{widened\_interval}: Parameterized interval expansion $\lambda \ge 1.0$ (GBR: \texttt{regime\_hcp}, $\lambda=1.00$; LSTM: \texttt{global\_hcp}, $\lambda=1.00$).
      \item \texttt{point\_safety\_margin}: Thresholding $p \le 1 + m$ (GBR: $m=5$; LSTM: $m=4$).
      \item \texttt{earlier\_intervention}: Elevated thresholding $L \le h$ (GBR: \texttt{regime\_hcp}, $h=1$; LSTM: \texttt{global\_hcp}, $h=1$).
    \end{enumerate}
  \item \textbf{Supported Logic}: For both GBR and LSTM, lower 98.75\% CI of $\Deltacost > 0.20$ SCU against \textbf{every single prespecified control configuration}, upper 98.75\% CI of marginal, regime, and stage coverage error $< 0.05$, and observed fleet failure rate increase $\le 0.0$ against all controls.
  \item \textbf{Unsupported Logic}: For either model, upper 98.75\% CI of $\Deltacost < 0.20$ SCU against any control configuration, lower 98.75\% CI of any coverage error $> 0.05$, or fleet failure rate increases over any control.
  \item \textbf{Inconclusive Logic}: No condition fails, but at least one decision interval covers its boundary, or inference is degenerate.
\end{itemize}

\subsubsection*{Research Question 4 (RQ4): Robustness and Generalization}
\begin{itemize}
  \item \textbf{Question}: \textit{Does the proposed method retain its calibration and maintenance-decision benefits across operating conditions, base RUL predictors, datasets, and repeated experimental runs?}
  \item \textbf{Null Hypothesis ($H_{0,4}$)}: The proposed method fails to achieve aggregate improvements ($\Deltacost > 0.20$ SCU, $\Deltaerror > 0.01$) across the evaluation matrix, or exhibits material harm ($\Deltacost < -0.20$ SCU or $\Deltaerror < -0.01$) in any individual configuration cell.
  \item \textbf{Alternative Hypothesis ($H_{1,4}$)}: The proposed method achieves aggregate improvements across the evaluation matrix without material harm in any individual configuration cell.
  \item \textbf{Evaluation Matrix}: Six mandatory cells: FD001 $\times$ GBR, FD001 $\times$ LSTM, FD002 $\times$ GBR, FD002 $\times$ LSTM, FD004 $\times$ GBR, FD004 $\times$ LSTM, weighted equally at $w_k = 1/6$ per cell.
  \item \textbf{Supported Logic}: Lower 98.75\% CI of weighted aggregate $\Deltacost > 0.20$ SCU, lower 98.75\% CI of weighted aggregate $\Deltaerror > 0.01$ for regimes and stages, lower 98.75\% CI of $\Deltacost > -0.20$ SCU in every cell, lower 98.75\% CI of $\Deltaerror > -0.01$ in every cell, and all cell safeguards hold across all 6 cells.
  \item \textbf{Unsupported Logic}: Upper 98.75\% CI of aggregate $\Deltacost < 0.20$ SCU or aggregate $\Deltaerror < 0.01$, or upper 98.75\% CI of cell $\Deltacost < -0.20$ SCU or $\Deltaerror < -0.01$ in any individual cell, or any cell safeguard fails.
  \item \textbf{Inconclusive Logic}: No condition fails, but at least one decision interval covers its boundary, or inference is degenerate.
\end{itemize}

\subsection*{Summary Evaluation Contract Matrix}

\begin{table}[h!]
\centering
\small
\begin{tabular}{p{0.8cm} p{1.8cm} p{2.8cm} p{2.8cm} p{2.0cm} p{3.2cm}}
\toprule
\textbf{ID} & \textbf{Experiment} & \textbf{Primary Outcome} & \textbf{Key Comparators} & \textbf{Statistical Unit} & \textbf{Practical Threshold (98.75\% CI)} \\
\midrule
\textbf{RQ1} & Experiment A & Maintenance cost per engine ($\Deltacost$) & \texttt{global\_hcp} \& \texttt{regime\_hcp} & Engine ($N=52$) & $\Deltacost > 0.20$ SCU per engine \\
\textbf{RQ2} & Experiment B & Worst-group coverage error ($\Deltaerror$) & \texttt{global\_hcp} \& \texttt{regime\_hcp} & Engine ($N=52$) & $\Deltaerror > 0.01$; $\Dwidth, \Dscore < 1.0$ \\
\textbf{RQ3} & Experiment C & Utility vs.\ matched controls ($\Deltacost$) & 4 named matched controls & Engine ($N=52$) & $\Deltacost > 0.20$ SCU vs.\ all controls \\
\textbf{RQ4} & Experiment D & Weighted aggregate \& cell non-harm & 6 cells ($w_k = 1/6$ each) & Engine ($N=122$) & Agg: $>0.20, >0.01$; Harm: $>-0.20, >-0.01$ \\
\bottomrule
\end{tabular}
\caption{IntelliTwin Confirmatory Hypotheses, Metrics, Comparators, and Frozen Thresholds (ADR-018).}
\end{table}
